\item La rapidez con la que una persona en reposo convierte energía alimenticia se llama \textit{tasa metabólica basal} (TMB). Suponga que la energía interna resultante deja el cuerpo de la persona por radiación y convección de aire seco. Cuando usted trota, la mayor parte de la energía alimenticia que quema por arriba de su TMB se convierte en energía interna que elevaría su temperatura corporal si no fuera eliminada. Suponga que la evaporación por transpiración es el mecanismo para eliminar esta energía. Suponga que una persona está trotando para la “máxima quema de grasa” convirtiendo energía alimenticia a la rapidez de $400\text{ kcal/h}$ por arriba de su TMB y expulsando energía por trabajo a la rapidez de 60.0 W. Suponga que el calor de evaporación del agua a la temperatura corporal es igual a su calor de vaporización a $100^\circ\text{C}$. (a) Determine la rapidez por hora a la cual el agua debe evaporarse de su piel. (b) Cuando usted metaboliza grasa, los átomos de hidrógeno en la molécula de grasa se transfieren al oxígeno para formar agua. Suponga que el metabolismo de 1.00 g de grasa genera 9.00 kcal de energía y produce 1.00 g de agua. ¿Qué fracción del agua que el corredor necesita es proporcionada por el metabolismo de grasas?
